\subsection{Bitwise Identities}

\subsubsection{Arithmetic Identities}

\textbf{Sum using AND and OR:}
\begin{equation}
    a + b = (a \& b) + (a | b)
\end{equation}

\textbf{Sum using XOR and AND:}
\begin{equation}
    a + b = (a \oplus b) + 2(a \& b)
\end{equation}

\textbf{XOR using AND and OR:}
\begin{equation}
    a \oplus b = \sim(a \& b) \& (a | b)
\end{equation}

\subsubsection{Bit Manipulation Operations}

Let $x$ denote the bit position (0-indexed from the right).

\textbf{Set (turn on) the $x$-th bit:}
\begin{equation}
    n \mid (1 \ll x)
\end{equation}

\textbf{Flip (toggle) the $x$-th bit:}
\begin{equation}
    n \oplus (1 \ll x)
\end{equation}

\textbf{Clear (turn off) the $x$-th bit:}
\begin{equation}
    n \& \sim(1 \ll x)
\end{equation}

\textbf{Check if the $x$-th bit is set:}
\begin{equation}
    n \& (1 \ll x) \neq 0
\end{equation}

\subsubsection{Useful Tricks}

\textbf{Clear all trailing ones:}
\begin{equation}
    n \& (n + 1)
\end{equation}

Example: $0011\ 0111 \rightarrow 0011\ 0000$

\textbf{Set the rightmost zero bit:}
\begin{equation}
    n | (n + 1)
\end{equation}

Example: $0011\ 0101 \rightarrow 0011\ 0111$

\textbf{Isolate the least significant bit (LSB):}
\begin{equation}
    n \& (-n)
\end{equation}

Example: $0011\ 0100 \rightarrow 0000\ 0100$

\textbf{Remove the least significant bit:}
\begin{equation}
    n \& (n - 1)
\end{equation}

Example: $0011\ 0100 \rightarrow 0011\ 0000$

\textbf{Check if $n$ is a power of 2:}
\begin{equation}
    (n \& (n - 1)) = 0 \quad \text{and} \quad n \neq 0
\end{equation}

Powers of 2 have exactly one bit set in their binary representation.
